\chapter{1926:The Svedberg}

\label{chap:chemistry1926}

The Nobel Prize in Chemistry for the year 1926 was awarded to The Svedberg.

\textbf{Motivation:}

for his work on disperse systems

\section{The Svedberg}

\subsection*{Early Life and Education}

Theodor Svedberg, known throughout his scientific career as The Svedberg, was born on 1884-08-30 in Valbo, Sweden.
He was a Swedish physical chemist who developed the ultracentrifuge and used it to measure the molecular weights of proteins and other macromolecules.

\subsection*{Career and Major Research}

Svedberg studied at the University of Uppsala, where he received his doctorate in 1908.
He was appointed professor of physical chemistry at Uppsala in 1912 and remained there for most of his career.
His early research concerned colloids, fine dispersions of particles in a liquid, and he developed methods for characterizing such disperse systems by measuring particle size and sedimentation.
Beginning in the early 1920s, Svedberg and his collaborator Herman Rinde constructed the ultracentrifuge, an instrument that spins liquid samples at very high speed so that dissolved molecules sediment at rates determined by their size.
With the ultracentrifuge, Svedberg showed that proteins are large molecules with definite, reproducible molecular weights, and he measured proteins such as haemoglobin, proving that molecules from different species have different molecular weights.

\subsection*{Nobel Prize-winning Work}

The work for which The Svedberg was awarded the Nobel Prize in Chemistry in 1926 was described by the Nobel Committee as: "for his work on disperse systems".

Svedberg's prize recognized his long series of investigations on colloids and other disperse systems.
Using the ultracentrifuge, he measured how rapidly colloidal particles and dissolved macromolecules sediment under centrifugal force, obtaining the size distribution of particles and, for proteins, the molecular weight of the individual molecules.
These experiments settled a long-running debate by demonstrating that proteins are well-defined molecules of characteristic size rather than indefinite aggregates.

\subsection*{Significance and Impact}

The ultracentrifuge gave chemists the first reliable tool for determining the molecular weights of proteins and other biological macromolecules, and it opened the way to a quantitative understanding of their behavior.
The unit of sedimentation rate, the svedberg (S), is named in his honour and remains standard in biochemistry and molecular biology.

\subsection*{Later Career and Honors}

The Svedberg continued his scientific work at Uppsala until his retirement in 1949.
He was a member of the Royal Swedish Academy of Sciences and held honorary degrees from several universities.
He died on 1971-02-25 in Sweden.

\subsection*{Legacy}

The legacy of The Svedberg endures through the lasting impact of his scientific contributions.
Ultracentrifugation remains a fundamental technique in the study of proteins and other macromolecules, and the Svedberg unit is used in laboratories throughout the world.

\subsection*{Selected Publications}

\begin{itemize}

\item Svedberg, T. (1912). "Die Existenz der Moleküle."

\item Svedberg, T., \& Rinde, H. (1924). "The Ultracentrifuge, a New Instrument for the Determination of Size and Distribution of Size of Particle in Amicroscopic Colloids." Journal of the American Chemical Society.

\end{itemize}

\clearpage

\section*{Chapter Summary}

The 1926 prize recognized The Svedberg for his work on disperse systems. His ultracentrifuge proved that proteins are large molecules of definite molecular weight and gave chemists a reliable method for characterizing colloids and macromolecules.
